% Generated draft: requires mathematical and source review.
\documentclass[11pt]{article}
\usepackage[margin=25mm]{geometry}
\usepackage{fontspec}
\setmainfont{Latin Modern Roman}
\usepackage{amsmath,amssymb,mathtools,unicode-math}
\setmathfont{Latin Modern Math}
\usepackage{enumitem,needspace,xurl,hyperref}
\setlength{\parindent}{0pt}
\setlength{\parskip}{5pt}
\setlength{\emergencystretch}{3em}
\newcommand{\N}{\mathbb{N}}
\newcommand{\Z}{\mathbb{Z}}
\newcommand{\Q}{\mathbb{Q}}
\newcommand{\R}{\mathbb{R}}
\newcommand{\C}{\mathbb{C}}
\newcommand{\F}{\mathbb{F}}
\newcommand{\A}{\mathbb{A}}
\newcommand{\PP}{\mathbb{P}}
\newcommand{\E}{\mathbb{E}}
\newcommand{\eps}{\varepsilon}
\newcommand{\dd}{\,\mathrm{d}}
\begin{document}
\section*{UnsolvedMath ID: 20000327}
\textbf{Source problem code:} catalogue entry
\subsection{Definitions and mathematical statement}
For an algebraic group $G$ acting on a quasi-projective scheme $U$, the quotient stack $[U/G]$ assigns to a test scheme $T$ the groupoid of $G$-torsors over $T$ with equivariant maps to $U$. The AIM problem asks whether every such quotient stack can be a moduli stack for some class of polarized manifolds. The original item does not specify affine versus arbitrary groups, ground field, polarization convention, or whether an equivalence with a full moduli substack is intended. The subsequent paragraph about interiors of canonically polarized moduli and period-domain analogies belongs to the next section, rather than adding hypotheses to this question.

For a linearly polarized manifold, interpret the data as a smooth projective variety $X$ with an ample line bundle $L$, not merely a numerical polarization class. Isomorphisms preserve $L$ up to isomorphism; scalar rigidification removes the scalar $\mathbb G_m$ acting on the line bundle. The supplied report notes that stabilizers in this convention are affine. If $A$ is a positive-dimensional abelian variety, then
\[
 BA=[\operatorname{Spec}k/A]
\]
is a quotient of a quasi-projective scheme but has non-affine stabilizer $A$. The report therefore refutes the unrestricted-group formulation for this line-bundle polarization convention. A different convention using only polarization classes requires separate treatment.

Its positive result is conditional on a fixed smooth connected projective parent $W$, a $G$-invariant quasi-projective $U\subset W$, and specified canonical-power, linearization, ampleness and full-automorphism hypotheses. For $u:T\to U$, the proposed family is the point blowup
\[
 W_u=\operatorname{Bl}_{\Gamma_u}(W\times T)
\]
with its chosen ample line bundle $L_u$. The report claims an isomorphism of relative fppf sheaves, functorially over arbitrary test schemes,
\[
 \operatorname{Trans}_G(u,v)\cong
 \operatorname{Isom}_{\mathrm{pol,rig}}((W_u,L_u),(W_v,L_v)),
 \qquad\operatorname{Trans}_G(u,v)(T')=
 \{g\in G(T'):g u_{T'}=v_{T'}\}.
\]
Thus $[U/G]$ maps fully faithfully to the scalar-rigidified polarized moduli stack and is identified with its full essential image. The claimed sheaf identification retains infinitesimal and nonreduced arrows, not just geometric-point automorphisms.

The summary does not actually list the canonical-power or automorphism hypotheses, the polarization formula or exponent, the parent dimension, or reconstruction proof. Accordingly this conditional special case is not a theorem realizing every affine-group quotient, nor a claim that the essential image is open or closed. In particular the point-blowup mechanism needs hypotheses ensuring the point can be recovered from the polarized blowup.
\subsection{Short English statement}
Ask whether quotient stacks can be realized as moduli of polarized smooth projective varieties. The report gives a non-affine-inertia obstruction for arbitrary groups, and a conditional construction using polarized point blowups for certain affine-group actions. Its positive statement concerns full relative isomorphism sheaves, but the essential construction hypotheses are missing from the supplied summary.
\subsection{Sources}
\begin{itemize}
\item Local registry: \url{data/unsolvedmath/problems.json}, permanent ID $20000327$, source code \texttt{catalogue entry}; original item, following-section context and supplied research summary.
\item AIM workshop ``Compact moduli spaces and birational geometry'', December 6--10, 2004, Problem 1.8: \url{https://aimath.org/WWN/birational/birational.pdf}. Canonical location \texttt{aim-algebraic-geometry-notes.json}, \texttt{notes[326]}; original remarks and literature field empty. The supplied verification mentions official TeX context but does not include its lead-in; no source or PDF was opened here.
\item Supplied machine-generated research attempt 1: partial result; problem partially solved at the run and in the registry; rigor metadata ``proved'', novelty confidence low. No authoritative citation or full hypotheses for the inertia and point-blowup claims are supplied. This non-peer-reviewed aid requires independent verification before citation.
\end{itemize}
\end{document}
